\section{Method}
\label{sec:method}

Given as input a monocular video sequence with dynamic objects,
our goal is to learn a neural network $f_\theta$ that can output a 4D representation of its geometry along with dense point-wise correspondences, \emph{simultaneously}:
\begin{equation}
    \label{eq:goal}
    f_\theta: \{\bm{I}_i\}_{i=1}^N \rightarrow
    \{\bm{X}_i, \bm{V}_{i\rightarrow i+1}\}_{i=1}^N.
\end{equation}
$\mathcal{I}=\{\bm{I}_i\}_{i=1}^N$ is the input monocular video sequence with $N$ frames,
where each frame
$\bm{I}_i\in\mathbb{R}^{H\times W\times 3}$
is an RGB image.
The network
$f_\theta$
predicts a viewpoint-invariant point map
$\bm{X}_i\in\mathbb{R}^{H\times W\times 3}$ for each frame $i$,
and the 3D scene flow\footnote{Unless otherwise noted,
we simplify $\bm{V}_{i\rightarrow i+1}$ as $\bm{V}_{i}$ for clarity.}
$\bm{V}_{i\rightarrow i+1}\in\mathbb{R}^{H\times W\times 3}$
between each pair of consecutive frames $i$ and $i+1$.
Both the point map and scene flow are represented in a shared \emph{world coordinate system}.
Note that,
since we only predict forward scene flow,
the last frame
$N$
does not have a corresponding flow prediction,
i.e., we do not supervise $\bm{V}_{N\rightarrow N+1}$.


To smoothly model the long-term motion and enable generalization to diverse scenes,
we build $f_\theta$ upon a pretrained video generator,
where $\boldsymbol{\theta}$ denotes the learnable parameters.
Our framework is illustrated in~\cref{fig:overview}.
We first introduce our unified 4D representation in~\Cref{sec:method_pre}.
Then,
in~\Cref{sec:method_vae},
we present a dedicated 4D\footnote{Here we use 4D VAE to refer to the fused geometry and motion VAEs.} VAE architecture that jointly encodes geometry and motion into a unified latent space.
Finally,
in~\Cref{sec:method_training},
we describe the overall training and inference strategy for our model.


\subsection{Unified Geometry \& Motion Representation}
\label{sec:method_pre}

Here,
like in DUSt3R~\cite{wang2024dust3r},
we define the point maps and scene flows in the coordinate system of the first frame,
which serves as the world coordinate system.
In particular,
the point map
$\bm{X}_i \in \mathbb{R}^{H\times W\times 3}$
stores the 3D coordinates
$(x, y, z)$
of each pixel from the frame
$i$
in the world coordinate system,
while the scene flow
$\bm{V}_{i} \in \mathbb{R}^{H\times W\times 3}$
represents the 3D motion vector \((\Delta x, \Delta y, \Delta z)\)
of each pixel from the frame
$i$
to
$i+1$.
Ideally, the \emph{deformed point map}
\begin{equation}
    \bm{X}_i^{d} = \bm{X}_i + \bm{V}_i
\end{equation}
should be spatially aligned with the point map of the next frame
$\bm{X}_{i+1}$. 
However,
due to viewpoint changes,
$\bm{X}_i^{d}$
and
$\bm{X}_{i+1}$
are not in one-to-one correspondence in pixel space, 
as they represent different frame contents,
as illustrated in~\cref{fig:representation}.
Note that,
% following the very recent work~\cite{sucar2025dynamic,st4rtrack2025},
our scene flow is also defined directly in the \emph{(world) coordinate system},
meaning that each
$\bm{V}_i$
represents the motion vector \((\Delta x, \Delta y, \Delta z)\)
in the world space,
naturally eliminating camera-induced motion components.

Such a unified geometry-motion representation offers several advantages:
%  for both modeling and learning:
1) \emph{Camera-free modeling.} 
Similar to DUSt3R~\cite{wang2024dust3r},
defining the geometry and motion in a chosen world coordinate system removes the need for additional camera pose estimation.
2) \emph{Temporal consistency.} 
In a continuous video sequence, geometry and motion are temporally coherent. 
Modeling them jointly in the same coordinate system makes them easier to learn.
3) \emph{Richer motion modeling.}
Unlike existing methods~\cite{sucar2025dynamic,st4rtrack2025},
we define scene flow between every pair of consecutive frames in the video,
rather than only between the first frame and others.
Consequently, this representation is less sensitive to occlusions induced by viewpoint variations and remains capable of capturing motion information of newly emerging dynamic objects in subsequent frames.


\begin{figure}
    \centering
    \includegraphics[width=\linewidth]{figs/data_representation_figure.pdf}
    \captionof{figure}{\textbf{Geometry and Motion representation.}
    % We define both point maps and scene flows in the world coordinate system.
    For a pixel $p_i$ in frame $\bm{I}_i$,
    $\bm{X}_i$ is its corresponding 3D point.
    As this 3D point moves,
    we use $\bm{X}_i^d$ to represent the moved point and $\bm{V}_{i} = (\Delta x, \Delta y, \Delta z)$ to represent the motion.
    Ideally,
    $\bm{X}_i^d$ should align with a matching point $\bm{X}_{i+1}$ in next frame $\bm{I}_{i+1}$.
    However, their pixel indexes are totally different ($p_i$ vs. $p_{i+1}$) and $p_{i+1}$ might even be out of view due to camera/object motion,
    making it impossible to build one-to-one correspondence between $\bm{X}_i^d$ and $\bm{X}_{i+1}$.
    }%
    \label{fig:representation}
    \vspace{-1em}

\end{figure}
\subsection{Unified 4D Geometry-Motion VAE}
\label{sec:method_vae}

Here, we describe how to effectively encode the above 4D representation into a latent space,
which can then be used as the target for a video generator.
Recent works on geometric diffusion models~\cite{ke2024repurposing,zhang2024world,jiang2025geo4d} encode only 3D geometric attributes, neglecting explicit motion modeling for dynamic scenes.
In contrast, we design a novel 4D VAE architecture that jointly encodes geometry and motion into a unified 4D latent,
as illustrated in~\cref{fig:overview}.

To leverage the priors of \emph{pretrained} video generators,
it is widely believed that
\emph{the input to the VAE should be strictly aligned with the original data distribution of the pre-trained diffusion model}~\cite{ke2024repurposing, zhang2024world, jiang2025geo4d}.
That is,
a naive approach is to directly rescale 3D attributes (such as disparity~\cite{ke2024repurposing} and point maps~\cite{zhang2024world}) into the range $[-1, 1]$,
by performing a max normalization,
and then encode them using the frozen VAE weights.
However,
the world-coordinate 3D attributes are typically unbounded, with coordinates spanning $(-\infty, +\infty)$,
in contrast to images with bounded pixel ranges $[0, 255]$.
Moreover, the distribution of 3D attributes is inherently distinct from that of natural RGB images.
Hence,
in this work,
we investigate a fundamental question:
\emph{Is strict alignment with the diffusion model's input space essential for finetuning diffusion models?}


\begin{figure}
    \centering
    \includegraphics[width=\linewidth]{figs/rescale_figure.pdf}
    \caption{\textbf{Results of different normalization and VAE training strategies}.
    For outdoor scenes with significant variations in depth (the second row), 
    the \emph{original VAE} fails to recover the scene structure.
    Even with decoder fine-tuning,
    the reconstruction quality remains poor.
    Our proposed \emph{mean} normalization and VAE training strategy significantly improve reconstruction quality.
    }
    \label{fig:rescale}
    \vspace{-1em}
\end{figure}


\paragraph{Geometry VAE with Revised Normalization}
To answer the above question,
one key insight of our model is a slightly adjusted point map normalization strategy for the Geometry VAE.
Note that,
unlike the max normalization to $[-1, 1]$ commonly used in existing geometric diffusion models~\cite{ke2024repurposing,zhang2024world,jiang2025geo4d},
we instead apply \emph{canonical normalization} to each sequence of world-coordinate point maps:
\begin{equation}
    \label{eq:normalization}
    \hat{\bm{X}}_i = \frac{\bm{X}_i-\mu}{S},
\end{equation}
where
$\mu=\frac{1}{|\mathcal{D}|}\sum_{d\in\mathcal{D}} \bm{X}_d$
is the mean of all valid points,
denoted by
$\mathcal{D}$,
in the point map sequence
$\cup_{i=1}^N\{\bm{X}_i\}$
and
$S=\frac{1}{|\mathcal{D}|}\sum_{d\in\mathcal{D}} \big\| \bm{X}_d - \mu \big\|_2 + \varepsilon$ 
is the mean distance for scale normalization with a small constant
$\varepsilon$
for numerical stability.
This normalization maintains the scale invariance of point maps,
while significantly improving the reconstruction quality of the Geometry VAE and better preserving finer structural details compared to max normalization~\cite{wang2024dust3r,wang2025vggt},
especially when handling large-scale outdoor scenes, as shown in~\cref{fig:rescale}.

Here, we finetune the entire encoder-decoder using our new normalization strategy, which provides a more flexible input distribution.
We define the training objective as:
\begin{equation}
\label{eq:geometry_loss}
\mathcal{L}_G = \mathcal{L}_{\text{point}} + \lambda_{d}\mathcal{L}_{\text{depth}} + \lambda_{n}\mathcal{L}_{\text{normal}},
\end{equation}
where
$\mathcal{L}_{\text{point}}$
is the MSE loss for point map reconstruction,
$\mathcal{L}_{\text{depth}}$
is a multi-scale loss computed on the projected depth maps,
and
$\mathcal{L}_{\text{normal}}$ enforces consistency of surface normals, following ~\cite{xu2025geometrycrafter, wang2024moge}.
The difference is that we are encoding world-based point clouds.
Hence, we normalize the ground-truth camera poses together with the point clouds so that
we can use scale-aligned camera parameters to project the point clouds into depth maps.
Experiments show that this supervision, similar to the multimodal fusion in Geo4D~\cite{jiang2025geo4d}, improves the reconstruction quality of point clouds.

Here we also tried using Kullback–Leibler (KL) divergence loss~\cite{kullback1951information} to constrain the distribution of the latent to a standard Gaussian distribution,
but found that it led to a significant drop in VAE performance.

The proposed normalization and VAE training strategy consistently improve the performance of the VAE and the downstream diffusion U-Net in~\cref{tab:ablation_geometry},
suggesting that \emph{strict alignment with the diffusion model's input and latent space is not always necessary},
especially for the 3D attributes. 

\paragraph{Motion VAE}
A simple way to model scene flow is to train a separate motion VAE with the same architecture as the Geometry VAE.
However, because motion and geometry are inherently correlated,
learning motion independently may be suboptimal.
We therefore explore several fusion strategies between geometry and motion:
(1) \textit{no fusion}, where geometry and motion are encoded separately without any interaction;
(2) \textit{offset fusion}, inspired by LayerDiffuse~\cite{zhang2024layerdiffuse},
where the motion latent is added as an offset to the geometry latent; and
(3) \textit{unified fusion}, where the geometry and motion latents are concatenated into a unified 4D latent
and passed to the motion VAE decoder to reconstruct the scene flow.
As reported in~\cref{tab:ablation_motion}, although the unified concatenation strategy does not yield the best reconstruction quality at the VAE stage,
it leads to superior performance in the subsequent diffusion U-Net.
During Motion VAE training, we freeze the Geometry VAE's parameters to preserve its learned geometric priors.
The training objective is formulated as:
% \begin{equation}
%     \label{eq:motion_loss}
%     \mathcal{L}_{M} = \mathcal{L}_{\text{sceneflow}} + \lambda_{\text{reg}}\mathcal{L}_{\text{reg}},
% \end{equation}
% \begin{equation}
%     \label{eq:motionreg_loss}
%     \mathcal{L}_{\text{reg}} = ,
% \end{equation}
\begin{equation}
\label{eq:motion_loss}
\mathcal{L}_{\text{M}} = 
\underbrace{
\frac{1}{|\mathcal{D}|} \sum_{d \in \mathcal{D}} \, \| \hat{\bm{V}}_d - \bm{V}_d \|_2^2
}_{\text{Scene flow reconstruction loss}}
+ \lambda_{\text{reg}} \underbrace{
\frac{1}{|\mathcal{N}|} \sum_{n \in \mathcal{N}} \, \| \hat{\bm{V}}_n \|_2^2
}_{\text{Zero-flow regularization}},
\end{equation}
where $\hat{\bm{V}}_d$ is the predicted scene flow, ${\bm{V}}_d$ is the ground truth,
$\mathcal{D}$ denotes valid pixels,
and $\mathcal{N}$ denotes all pixels.
The first term denotes the MSE loss on the valid scene flow,
and the second term is a regularization term to encourage the scene flow to zero,
following the as-static-as-possible assumption. 


By combining the Geometry VAE and Motion VAE into a unified 4D VAE, 
we successfully achieve an integrated representation of geometry and motion within a single latent space, 
enabling efficient 4D scene encoding and decoding.

\subsection{Model Training}
\label{sec:method_training}

\paragraph{Training Data}
Dynamic datasets with annotated 3D geometry and dense scene flow are difficult to collect in real-world settings.
Therefore,
we rely on synthetic datasets for training the scene flow estimation task.
In particular,
we divide our training data into two categories:
(1) \emph{Geometry Datasets:}
Dynamic Replica~\cite{karaev2023dynamicstereo},
GTA-SFM~\cite{wang2020gtasfm},
MatrixCity~\cite{li2023matrixcity},
MVS-Synth~\cite{huang2018deepmvs},
Point Odyssey~\cite{zheng2023pointodyssey},
TartanAir~\cite{wang2020tartanair},
ScanNet++~\cite{yeshwanth2023scannet++},
BlinkVision~\cite{li2024blinkvision},
OmniWorld~\cite{zhou2025omniworld}
and Synthia~\cite{ros2016synthia};
(2) \emph{Geometry-and-Motion Datasets:}
Kubric~\cite{greff2021kubric},
Spring~\cite{Mehl2023_Spring},
and Virtual KITTI 2~\cite{cabon2020virtual}.
The first category provides only geometric data,
including per-frame depth maps,
camera intrinsics and extrinsics.
Following DUSt3R~\cite{wang2024dust3r},
we express the ground-truth point clouds into a shared first frame coordinate system.
The second category additionally provides dense scene flow annotations.
During \textit{geometry reconstruction} training,
we use both dataset groups (1)+(2),
whereas during \textit{motion reconstruction} training,
only datasets in group (2) are employed.

\paragraph{Training Strategy}
We adopt a two-stage training pipeline for the VAE components.
We begin by training the Geometry VAE independently to capture scene geometry.
Next,
we train the Motion VAE while keeping the Geometry VAE frozen,
thereby preserving its learned geometric priors.
After convergence,
we combine them into a unified 4D VAE,
whose parameters remain frozen during the training of the diffusion U-Net.
For U-Net training, we combine the datasets from groups (1) and (2) to provide geometry supervision, and use only the datasets from group (2) for motion supervision.
Following prior works~\cite{hu2024-DepthCrafter,xu2025geometrycrafter} in employing EDM~\cite{karras2022elucidating} pre-conditioning,
our framework supports both the \textit{deterministic} and \textit{denoising} paradigms.
For the deterministic paradigm, the training objective is defined as:
\begin{equation}
\mathcal{L}_{\text{deterministic}} = \mathcal{L}_{\text{latent}} + \lambda_G \mathcal{L}_G + \lambda_M \mathcal{L}_{\text{M}},
\end{equation}
where $\mathcal{L}_G$ is the geometry reconstruction loss defined in~\eqref{eq:geometry_loss},
$\mathcal{L}_{M}$ is the motion reconstruction loss defined in~\eqref{eq:motion_loss}, and $\mathcal{L}_{\text{latent}}$ denotes the latent-space diffusion loss:
\begin{equation}
\mathcal{L}_{\text{latent}} =
\underbrace{
\frac{1}{N} \sum_{N} 
\| 
\hat{\mathbf{z}}^{\text{G}}_i - \mathbf{z}^{\text{G}}_i
\|_2^2
}_{\text{geometry latent supervision}}
+
\underbrace{
\frac{1}{N-1} \sum_{N-1}
\| 
\hat{\mathbf{z}}^{\text{M}}_i - \mathbf{z}^{\text{M}}_i
\|_2^2
}_{\text{motion latent supervision}},
\end{equation}
where $N$ denotes the number of frames,
and $\hat{\mathbf{z}}^{\text{G}}_i,\mathbf{z}^{\text{G}}_i,\hat{\mathbf{z}}^{\text{M}}_i, \mathbf{z}^{\text{M}}_i$ are the denoised latent and the original latent of geometry and motion, respectively.
We only perform forward scene flow estimation,
thus discarding the motion latent in the last frame.
For the denoising paradigm, the objective simplifies to latent supervision:
% \begin{equation}
$\mathcal{L}_{\text{denoise}} = \mathcal{L}_{\text{latent}}$.
% \end{equation}
We found experimentally that the deterministic paradigm generally performs better, so we use it by default. Ablation experiments can be found in the supplementary materials.

This progressive and modular training pipeline allows the model to first acquire strong geometry and motion priors before integrating temporal reasoning,
ultimately enabling robust and coherent dense 4D reconstruction.

\paragraph{Implementation Details}
To inherit the strong priors from the video generator,
both the VAE and U-Net of our \method{} are initialized with the pretrained weights of SVD~\cite{blattmann2023stable},
and trained using the AdamW optimizer~\cite{loshchilov2017decoupled} with a learning rate of $1e$-$4$. 
We first train the Geometry VAE for 40{,}000 iterations,
followed by training the Motion VAE for 20{,}000 iterations. 
Subsequently,
we merge the Geometry VAE and Motion VAE into a unified 4D VAE,
and train the U-Net for another 40{,}000 iterations with the encoded 4D latent representations.
The batch size is set to 8 for VAE training and 25 for U-Net training.
All experiments are conducted on 8 GPUs with 40 GB of memory each and take about 3 days.
More implementation details are provided in the supplementary material.